% UDL UK Patent Application — Specification
% Title: Method for Quantifying Anomalies Using Multi-Spectrum Tensor
%        Decomposition Across Independent Mathematical Law Domains
% Applicant: Odeyemi Olusegun Israel
% Date: March 2026

\documentclass[12pt,a4paper]{article}
\usepackage[margin=2.5cm]{geometry}
\usepackage{amsmath,amssymb}
\usepackage{graphicx}
\usepackage{enumitem}
\usepackage{hyperref}
\usepackage{titlesec}
\usepackage{fancyhdr}
\usepackage{booktabs}
\setlength{\headheight}{14.5pt}

\pagestyle{fancy}
\fancyhf{}
\rhead{Patent Application}
\lhead{UDL --- Odeyemi O.I.}
\rfoot{Page \thepage}

\titleformat{\section}{\large\bfseries}{\thesection.}{0.5em}{}
\titleformat{\subsection}{\normalsize\bfseries}{\thesubsection.}{0.5em}{}

\begin{document}

\begin{center}
{\LARGE\bfseries PATENT SPECIFICATION}\\[1em]
{\Large Method for Quantifying Anomalies Using Multi-Spectrum\\
Tensor Decomposition Across Independent Mathematical Law Domains}\\[2em]
{\large Applicant: Odeyemi Olusegun Israel}\\[0.5em]
{\large Inventor: Odeyemi Olusegun Israel}\\[0.5em]
{\large Date of Filing: 4 March 2026}\\[0.5em]
{\large Application Number: [TO BE ASSIGNED]}\\[0.3em]
{\normalsize \textit{Filed pursuant to the Patents Act 1977}}
\end{center}

\vspace{2em}

% ============================================================
\section{TECHNICAL FIELD}
% ============================================================

The present invention relates to the field of anomaly detection and 
machine learning, and more particularly to a model-independent framework 
for detecting and characterising anomalies by measuring deviations across 
multiple independent mathematical representations of data.

% ============================================================
\section{BACKGROUND OF THE INVENTION}
% ============================================================

\subsection{Technical Problem}

Anomaly detection---the identification of observations that deviate 
significantly from expected patterns---is a foundational problem in data 
analysis with applications across finance, cybersecurity, industrial 
monitoring, medical diagnosis, and scientific research. Existing approaches 
suffer from several structural limitations:

\begin{enumerate}[label=(\alph*)]
\item \textbf{Single-perspective detection:} Standard anomaly detectors 
(Isolation Forest, Local Outlier Factor, One-Class SVM, autoencoders) 
produce a single anomaly score. They identify \emph{that} an observation 
is anomalous but not \emph{which mathematical aspect} makes it anomalous. 
This limits interpretability and actionability.

\item \textbf{Detector-dependent representations:} Each detector imposes its 
own implicit representation of ``normality.'' An observation that appears 
normal under one detector's representation (e.g., distance-based) may 
simultaneously violate a different mathematical law (e.g., spectral or 
dynamical). No single detector spans the full space of mathematical anomaly.

\item \textbf{Lack of attribution:} When an anomaly is detected, 
practitioners cannot determine whether the deviation is statistical, 
dynamical, spectral, geometric, or some combination thereof, without running 
multiple separate analyses.

\item \textbf{Coverage gaps:} High AUC (area under the ROC curve) does not 
guarantee high coverage---the anomaly scores may concentrate in a narrow 
band, failing to discriminate among different types of anomalies (the 
``coverage--AUC gap'' problem).

\item \textbf{Hyperparameter sensitivity:} Most detectors require per-dataset 
tuning of hyperparameters (contamination rate, number of neighbours, kernel 
bandwidth), limiting practical deployability.
\end{enumerate}

\subsection{Prior Art Limitations}

Known prior art includes Isolation Forest (Liu et al., 2008), Local Outlier 
Factor (Breunig et al., 2000), One-Class SVM (Sch\"olkopf et al., 1999), 
deep autoencoders, out-of-distribution detection methods (ODIN, energy-based 
methods), and ensemble methods that combine multiple instances of the same 
detector type. None of these provides multi-law-domain decomposition of 
anomaly type, per-law attribution, or a formal guarantee that anomalies 
cannot simultaneously appear normal across all mathematical perspectives.

% ============================================================
\section{SUMMARY OF THE INVENTION}
% ============================================================

The present invention provides a method and system for detecting and 
characterising anomalies by simultaneously evaluating observations under 
multiple independent mathematical law domains. The key contributions are:

\begin{enumerate}[label=(\roman*)]
\item The \textbf{Universal Deviation Principle}: a formal mathematical 
statement, supported by theorems, that if a family of operators 
$\{\Phi_k\}_{k=1}^{K}$ is \emph{deviation-separating} (i.e., the stacked 
Jacobian matrix has full rank), then no anomalous observation can 
simultaneously match normal data across all operators.

\item A \textbf{multi-spectrum anomaly tensor} $T_{i,k,d}$ constructed by 
projecting each observation through at least five independent mathematical 
law domains.

\item A \textbf{Magnitude--Direction--Novelty (MDN) decomposition} of the 
anomaly tensor that separates how much, what type, and how novel each 
deviation is.

\item A \textbf{Boundary-Centred Dimension Magnifier} that adaptively 
rescales tensor dimensions based on empirical separability.

\item A classification stage (QDA-Magnified or Hybrid auto-selection) that 
produces per-law-domain attributed anomaly scores.
\end{enumerate}

% ============================================================
\section{DETAILED DESCRIPTION OF THE INVENTION}
% ============================================================

\subsection{Universal Deviation Principle}

The method is founded on the following formal principle:

\textbf{Definition (Deviation-Separating Operator Family):} An operator 
family $\{\Phi_k\}_{k=1}^{K}$ is \emph{deviation-separating} if the stacked 
Jacobian matrix:
\begin{equation}
J(x) = \begin{bmatrix} D\Phi_1(x) \\ D\Phi_2(x) \\ \vdots \\ 
D\Phi_K(x) \end{bmatrix}
\end{equation}
has full column rank for all $x$ in the data domain.

\textbf{Theorem 1 (Local Deviation Guarantee):} If the operator family 
$\{\Phi_k\}$ is deviation-separating, then by the inverse function theorem, 
for any anomalous point $x^*$ not identical to the normal centroid 
$\mu_{\text{normal}}$, there exists at least one operator $\Phi_k$ under 
which $x^*$ deviates from the projection of normal data.

\textbf{Theorem 2 (Global Deviation Guarantee):} With a compact anomaly set, 
a quantitative lower bound on the minimum deviation across all operators can 
be established.

\textbf{Theorem 3 (Finite-Sample Guarantee):} Given $n$ normal training 
samples, the detection guarantee holds with probability at least 
$1 - \delta$ for a sample complexity depending on the operator family's 
covering number.

\textbf{Theorem 4 (Information-Optimal Fusion):} The optimal weights for 
combining operator outputs can be computed as the solution to a semidefinite 
programme (SDP).

\textbf{Theorem 5 (Greedy Operator Selection):} A greedy algorithm for 
selecting the most informative subset of operators from a large candidate 
set achieves a $(1 - 1/e)$ approximation to the optimal subset.

\subsection{Spectrum Operators (Law Domains)}

The method employs at least five fundamental mathematical law domains, each 
implemented as a spectrum operator $\Phi_k$:

\subsubsection{Statistical Operator ($\Phi_{\text{stat}}$)}

Measures deviations in distributional properties:
\begin{itemize}
\item Shannon entropy of the observation's local neighbourhood
\item Kullback--Leibler divergence from the reference normal distribution
\item Hellinger distance from the reference distribution
\item Skewness and kurtosis of local feature distributions
\end{itemize}

\subsubsection{Dynamical (Chaos) Operator ($\Phi_{\text{chaos}}$)}

Measures deviations in temporal dynamics:
\begin{itemize}
\item Lyapunov exponent estimated from the observation's temporal context
\item Recurrence rate from recurrence plot analysis
\item Approximate entropy measuring regularity of temporal patterns
\end{itemize}

\subsubsection{Spectral (Frequency) Operator ($\Phi_{\text{freq}}$)}

Measures deviations in frequency-domain representations:
\begin{itemize}
\item Power spectral density divergence from reference PSD
\item Frequency ratio comparing low-frequency to high-frequency energy
\item Spectral entropy measuring uniformity of the frequency spectrum
\item Spectral centroid shift from the reference centroid
\end{itemize}

\subsubsection{Geometric Operator ($\Phi_{\text{geom}}$)}

Measures deviations in geometric relationships:
\begin{itemize}
\item Mahalanobis distance from the reference centroid
\item Cosine similarity to the reference direction
\item Norm ratio comparing observation magnitude to reference magnitude
\end{itemize}

\subsubsection{Exponential Operator ($\Phi_{\text{exp}}$)}

Measures deviations under exponential amplification:
\begin{equation}
\Phi_{\text{exp}}(x) = \exp\left(\alpha \cdot \frac{x - \mu}{\sigma}\right)
\end{equation}
where $\mu$ and $\sigma$ are the reference mean and standard deviation and 
$\alpha$ is a scaling parameter. This operator amplifies small deviations 
that may be invisible under linear representations.

\subsubsection{Extended Operators}

In an enhanced embodiment, additional operators are provided:
\begin{itemize}
\item Phase operator: phase-space embedding analysis
\item Topological operator: persistent homology features
\item Kernel operator: kernel-space deviation
\item Rank operator: rank-based deviation statistics
\item Wavelet operator: multi-scale wavelet decomposition
\end{itemize}

In the preferred embodiment, 14 composable operators are available. An 
operator correlation audit confirms that only 1 of 91 operator pairs 
exhibits redundancy (Wavelet--Phase, $\rho = 0.92$); all others are 
sufficiently independent ($\rho < 0.85$).

\subsection{Anomaly Tensor Construction}

For each observation $x_i$, each operator $\Phi_k$ produces a vector of 
deviation measurements across $D$ dimensions. The full anomaly tensor is:
\begin{equation}
T_{i,k,d} \quad \text{for observation } i, \text{ operator } k, 
\text{ dimension } d
\end{equation}

This tensor captures the multi-law deviation profile of each observation.

\subsection{Magnitude--Direction--Novelty (MDN) Decomposition}

The anomaly tensor for each observation is decomposed into three 
interpretable components:

\begin{enumerate}[label=(\alph*)]
\item \textbf{Magnitude} $r_i = \|T_i\|$: a scalar measuring the total 
amount of deviation across all operators and dimensions.

\item \textbf{Direction} $\hat{D}_i = T_i / \|T_i\|$: a unit tensor 
indicating the type of deviation---which operators and dimensions contribute 
most.

\item \textbf{Novelty} $\theta_i$: the angular distance of $\hat{D}_i$ from 
reference deviation directions observed in the training data, measuring how 
novel the deviation pattern is.
\end{enumerate}

This decomposition enables three-axis interpretation: ``how much,'' ``what 
type,'' and ``how novel'' the anomaly is.

\subsection{Boundary-Centred Dimension Magnifier}

A key technical contribution is the Boundary-Centred Dimension Magnifier, 
which rescales each dimension of the anomaly tensor according to its 
empirical separability between normal and anomalous classes:

\begin{equation}
r'_d = \tanh\left(k_d \cdot \frac{r_d - b_d}{\sigma_d}\right)
\end{equation}

where:
\begin{itemize}
\item $r_d$ is the raw value in dimension $d$;
\item $b_d$ is the decision boundary location in dimension $d$;
\item $\sigma_d$ is the scale of the boundary region;
\item $k_d = \gamma \cdot \delta_d$ is the steepness, with $\gamma$ a 
global parameter and $\delta_d$ a per-dimension discriminability score.
\end{itemize}

The per-dimension discriminability $\delta_d$ is computed from three 
complementary measures:
\begin{enumerate}
\item Cohen's $d$ (effect size) between normal and anomalous distributions 
in dimension $d$;
\item An overlap metric measuring the overlap area of the two class 
distributions; and
\item Per-dimension AUC measuring the classification performance using 
dimension $d$ alone.
\end{enumerate}

In the preferred embodiment, the global parameter $\gamma = 5.0$ and the 
QDA regularisation parameter $\lambda = 10^{-4}$ are used as a single 
hyperparameter set across all datasets.

The magnifier has the effect of:
\begin{itemize}
\item Amplifying dimensions where normal and anomalous classes are 
well-separated (high $\delta_d$);
\item Suppressing dimensions where the classes overlap heavily (low 
$\delta_d$);
\item Centring the rescaling around the empirical decision boundary rather 
than the class centroids.
\end{itemize}

\subsection{Classification Stage}

The magnified anomaly tensor is classified using one of four strategies, 
selected automatically via cross-validation (Hybrid auto-selection):

\begin{enumerate}[label=(\alph*)]
\item \textbf{Fisher projection:} Linear discriminant projection followed by 
thresholding.
\item \textbf{QDA:} Quadratic discriminant analysis with regularised 
per-class covariance.
\item \textbf{QDA-Magnified:} QDA applied to the magnifier-rescaled tensor 
(preferred strategy).
\item \textbf{RankFusion:} A fully unsupervised alternative using rank 
statistics of per-operator deviations.
\end{enumerate}

The Hybrid auto-selection evaluates all four strategies on cross-validation 
folds and selects the strategy maximising coverage (the proportion of 
anomalies detected above threshold) while maintaining AUC above a minimum 
floor.

\subsection{Centroid Estimation}

Seven centroid estimation strategies are provided:
\begin{enumerate}
\item Sample mean
\item Trimmed mean
\item Geometric median
\item Minimum covariance determinant (MCD)
\item Medoid
\item Huber M-estimator
\item Density-weighted centroid
\end{enumerate}

Auto-selection chooses the centroid strategy that minimises within-class 
scatter of normal data in the magnified space.

\subsection{Coverage Metric}

The method introduces a coverage metric that complements AUC. Coverage 
measures the concentration of anomaly scores in the top-$k$ ranked 
observations:
\begin{equation}
\text{Coverage@}k = \frac{|\{x : \text{rank}(x) \leq k \text{ and } 
x \in \text{anomaly}\}|}{|\text{anomaly}|}
\end{equation}

The coverage metric addresses the AUC--coverage gap: detectors may achieve 
high AUC while concentrating all anomaly scores in a narrow band, failing 
to discriminate among different anomaly subtypes. In the preferred 
embodiment, mean coverage exceeds 91\%.

\subsection{Performance Characteristics}

Using a single hyperparameter set ($\gamma = 5.0$, $\lambda = 10^{-4}$):

\begin{center}
\begin{tabular}{lcc}
\toprule
\textbf{Dataset} & \textbf{QDA-Magnified AUC} & \textbf{Hybrid AUC} \\
\midrule
Synthetic & 1.000 & 0.999 \\
Mimic (camouflage adversarial) & 1.000 & 0.999 \\
Mammography & 0.947 & 0.911 \\
Shuttle & 0.989 & 0.991 \\
Pendigits & 0.777 & 0.960 \\
\midrule
\textbf{Mean} & \textbf{0.943} & \textbf{0.972} \\
\bottomrule
\end{tabular}
\end{center}

% ============================================================
\section{CLAIMS}
% ============================================================

\noindent\textit{The following claims define the scope of protection
sought. Where a claim refers to another claim, the dependency is by claim
number. Features of any dependent claim may be combined with any
independent claim or with other dependent claims.}
\vspace{0.8em}

\begin{enumerate}[label=\textbf{\arabic*.}]

% --- Independent Method Claim ---
\item A computer-implemented method for detecting anomalies in data, the 
method comprising the steps of:
\begin{enumerate}[label=(\alph*)]
\item receiving an input observation comprising a feature vector;

\item projecting the observation through a plurality of spectrum operators, 
each spectrum operator computing deviation measurements under a different 
mathematical law domain, the plurality comprising at least:
\begin{enumerate}[label=(\roman*)]
\item a statistical operator measuring distributional deviations;
\item a dynamical operator measuring temporal dynamical deviations;
\item a spectral operator measuring frequency-domain deviations;
\item a geometric operator measuring geometric relationship deviations; and
\item an exponential operator measuring deviations under exponential 
amplification;
\end{enumerate}

\item constructing an anomaly tensor from the deviation measurements across 
all spectrum operators;

\item decomposing the anomaly tensor into magnitude, direction, and novelty 
components;

\item applying a boundary-centred dimension magnifier that rescales each 
dimension of the anomaly tensor according to an empirically computed 
per-dimension discriminability score; and

\item classifying the magnified anomaly tensor to produce an anomaly score 
and per-law-domain attribution indicating which mathematical law domains 
contribute to the anomaly determination.
\end{enumerate}

% --- Independent System Claim ---
\item A system for detecting and characterising anomalies across multiple 
mathematical law domains, the system comprising:
\begin{enumerate}[label=(\alph*)]
\item a plurality of spectrum operator modules, each configured to compute 
deviation measurements of an input observation under a different mathematical 
law domain;

\item a tensor construction module configured to assemble the deviation 
measurements into a multi-dimensional anomaly tensor;

\item a decomposition module configured to decompose the anomaly tensor 
into magnitude, direction, and novelty components;

\item a magnifier module configured to rescale each dimension of the 
anomaly tensor using a boundary-centred rescaling function with 
per-dimension steepness computed from empirical class separability; and

\item a classifier module configured to produce an anomaly score and 
per-law-domain attribution from the magnified tensor;
\end{enumerate}
wherein the plurality of spectrum operators collectively satisfy a 
deviation-separating condition such that the stacked Jacobian matrix of the 
operators has full column rank.

% --- Dependent Claims ---
\item A method according to claim 1, wherein the boundary-centred dimension 
magnifier applies the rescaling function:
\[
r'_d = \tanh\left(k_d \cdot \frac{r_d - b_d}{\sigma_d}\right)
\]
where $k_d = \gamma \cdot \delta_d$, $\gamma$ is a global steepness 
parameter, and $\delta_d$ is a per-dimension discriminability score computed 
from Cohen's $d$, class distribution overlap, and per-dimension AUC.

\item A method according to claim 1, wherein the decomposition into 
magnitude, direction, and novelty comprises:
\begin{enumerate}[label=(\alph*)]
\item computing the magnitude as the norm of the anomaly tensor;
\item computing the direction as the normalised anomaly tensor; and
\item computing the novelty as the angular distance of the direction from 
reference deviation directions in the training data.
\end{enumerate}

\item A method according to claim 1, wherein the classifying step uses 
quadratic discriminant analysis with regularised per-class covariance 
matrices applied to the magnified anomaly tensor (QDA-Magnified).

\item A method according to claim 1, further comprising a hybrid 
auto-selection step that evaluates a plurality of classification strategies 
on cross-validation folds and selects the strategy maximising a coverage 
metric while maintaining AUC above a minimum threshold.

\item A method or system according to any preceding claim, further 
comprising an operator correlation audit step that verifies independence 
of the spectrum operators by confirming that pairwise correlation 
coefficients are below a redundancy threshold.

\item A method or system according to any preceding claim, further 
comprising a centroid auto-selection module configured to select from a 
plurality of centroid estimation strategies the strategy that minimises 
within-class scatter of normal data in the magnified space.

\item A method or system according to any preceding claim, wherein a 
single global hyperparameter set is used across all datasets without 
per-dataset tuning.

\item A method or system according to any preceding claim, further 
comprising the step of computing a coverage metric that measures the 
proportion of anomalies detected in the top-ranked observations, the 
coverage metric complementing AUC by quantifying score concentration.

\item A method or system according to any preceding claim, wherein the per-law-domain attribution identifies which mathematical law domains---among 
statistical, dynamical, spectral, geometric, and exponential---contribute 
most to each anomaly determination, enabling interpretable anomaly 
diagnosis.

\item A method according to claim 1, further comprising a greedy operator 
selection step that selects a subset of operators from a larger candidate 
set, the selection achieving a $(1 - 1/e)$ approximation to the optimal 
subset under a submodular information criterion.

\end{enumerate}

% ============================================================
\section{DRAWINGS}
% ============================================================

\textit{[To be filed separately: Diagram of the multi-spectrum operator 
architecture; illustration of the anomaly tensor construction; 
visualisation of the MDN decomposition; schematic of the boundary-centred 
dimension magnifier; operator independence correlation matrix; 
cross-dataset performance comparison.]}

% ============================================================
\newpage
\section*{ABSTRACT}
% ============================================================

A computer-implemented anomaly detection method projects each observation
through a plurality of independent mathematical law domains comprising
statistical, dynamical, spectral, geometric and exponential operators,
constructing a multi-dimensional anomaly tensor. The method is founded on
the principle that no anomaly can simultaneously match normal data across all
mathematically independent law domains. The anomaly tensor is decomposed into
magnitude, direction and novelty components. A boundary-centred dimension
magnifier rescales each tensor dimension according to its empirical
discriminative power, amplifying well-separated dimensions and suppressing
overlapping dimensions. A quadratic discriminant analysis classifier operating
on the magnified tensor produces anomaly scores with per-law-domain
attribution, enabling interpretable identification of which mathematical
aspect renders each observation anomalous. The method achieves a mean
ROC-AUC of 0.943 across synthetic, financial, aerospace and medical datasets
using a single hyperparameter set.

\end{document}
